% 第3章 阿秒电子脉冲的产生与表征方法
\chapter{阿秒电子脉冲的产生与表征方法}
\label{chap:method}

本章介绍实验系统的搭建与关键参数的表征，包括透射电子显微镜的光学集成方案、连续波激光驱动的阿秒电子脉冲产生、能量滤波器的优化配置、以及时间分辨率的标定方法。

\section{实验系统概述}
\label{sec:exp-overview}

% 实验系统整体介绍：TEM + CW激光 + 能量滤波器

\begin{figure}[htbp]
    \centering
    % TODO: 从素材库提取实验系统示意图
    \includegraphics[width=0.9\linewidth]{ch3/fig_exp_setup}
    \caption{\label{fig:ch3-setup} 阿秒电子显微术实验系统示意图。连续波激光（红色光路）同时照射压缩膜和样品。电子束（蓝色箭头）穿越压缩膜获得能量调制，在样品平面处形成阿秒脉冲串。能量滤波器选择特定能量侧带的电子成像。}
\end{figure}

\section{连续波激光驱动的电子脉冲压缩}
\label{sec:cw-compression}

% 3.2 CW激光参数选择
% 3.2.1 压缩膜（Si3N4）的制备与集成
% 3.2.2 激光-电子对准方案
% 3.2.3 调制深度|g|的测量与优化

\section{能量滤波器的对比度增强}
\label{sec:energy-filter}

% 3.3 能量滤波器配置
% 3.3.1 GIF（Gatan Imaging Filter）的参数设置
% 3.3.2 能量窗口宽度的选择
% 3.3.3 背景扣除与信噪比优化

\section{时间扫描与延迟控制}
\label{sec:delay-control}

% 3.4 时间延迟的精密控制
% 3.4.1 压电平移台的校准
% 3.4.2 时间零点的确定
% 3.4.3 相位锁定与稳定性

\section{时间分辨率的标定}
\label{sec:resolution-calibration}

% 3.5 时间分辨率的实验标定
% 3.5.1 从能量谱推算脉冲宽度
% 3.5.2 空间分辨率的测量（PSF法）
% 3.5.3 系统误差分析

\begin{table}[htbp]
    \caption{\label{tab:ch3-params} 实验系统关键参数}
    \begin{tabularx}{\linewidth}{c|X<{\centering}}
        \hline
        参数 & 数值 \\ \hline
        电子能量 $E_0$ & \SI{183}{keV} \\ \hline
        激光波长 $\lambda$ & \SI{785(5)}{nm} \\ \hline
        激光功率 $P$ & \SI{300(10)}{mW} \\ \hline
        脉冲宽度（FWHM）& \SI{800(50)}{as} \\ \hline
        空间分辨率 & \SI{50(10)}{nm} \\ \hline
        时间扫描步长 & $\sim$\SI{3}{as} \\ \hline
        能量滤波窗口 & $\Delta E > $\SI{1}{eV} \\ \hline
    \end{tabularx}
\end{table}

\section{本章小结}
\label{sec:method-summary}

本章介绍了阿秒电子显微术实验系统的搭建与关键参数表征。系统以连续波激光驱动，在常规TEM中实现了$\sim$\SI{800}{as}的时间分辨率和$\sim$\SI{50}{nm}的空间分辨率，为后续的纳米光子系统亚周期动力学研究提供了可靠的实验基础。
